%% kapa: exam=1 round=32 year=2021 period=1 form=A subject=civil_law subject_ko=민법 seq=39 number=39
%% source: Q-Net 감정평가사 기출문제 (cst00309) artlSeq=5208397 file=extracted/1008/5208397/A형/2021년도 제32회 감정평가사 1차 1교시 A형.pdf
%% conversion: kordoc 4.24.51 -> tools/gen1.py (auto); review=proofread-sonnet; flags=none
\begin{question}{39}
甲은 乙에 대한 2억 원의 채권을 담보하기 위하여 乙 소유 X토지와 Y건물에 대하여 각각 1번 공동저당권을 취득하였다. 그 후 丙은 乙에 대한 1억 6천만 원의 채권을 담보하기 위하여 X에 대하여 2번 저당권을, 丁은 乙에 대한 7천만 원의 채권을 담보하기 위하여 Y에 대하여 2번 저당권을 취득하였다. 그 후 丙이 경매를 신청하여 X가 3억 원에 매각되어 배당이 완료되었고, 다시 丁이 경매를 신청하여 Y가 1억 원에 매각되었다. 丁이 Y의 매각대금에서 배당받을 수 있는 금액은? (단, 경매비용·이자 등은 고려하지 않으며, 다툼이 있으면 판례에 따름)
\begin{choices}
  \item 0원
  \item 3,500만원
  \item 4,000만원
  \item 5,000만원
  \item 7,000만원
\end{choices}
\end{question}
\answer{4}
